Auch wenn sich diese Arbeit auf die Entwicklung einer barrierefreien Webanwendung im Sinne der Programmierung konzentriert, spielt die Planung und Konzeption eine entscheidende Rolle. Bereits in dieser frühen Phase müssen wichtige Aspekte der Barrierefreiheit berücksichtigt werden. Im Folgenden wird daher beschrieben, wie die Anwendung geplant und gestaltet wurde.
Der Prototyp mit dem Design lässt sich hier abrufen: 
\href{https://www.figma.com/proto/jnGNE7gzjqBn1uD3rOXifO/bachelor-thesis?node-id=1-79\&p=f\&t=NlQQAOGGDPRokds2-0\&scaling=min-zoom\&content-scaling=fixed\&page-id=0\%3A1\&starting-point-node-id=1\%3A79}{Design Prototyp}\footnote{\url{https://www.figma.com/proto/jnGNE7gzjqBn1uD3rOXifO/bachelor-thesis?node-id=1-79\&p=f\&t=NlQQAOGGDPRokds2-0\&scaling=min-zoom\&content-scaling=fixed\&page-id=0\%3A1\&starting-point-node-id=1\%3A79}}
\section{Vorgaben durch Mrija Manager}
Die Arbeit basiert auf dem Projekt Mrija Manager: Daher gibt es bei dem Konzept und bei der Gestaltung bestimmte Vorgaben.

\section{Seitenstruktur}
Die Konzeption von Mrija Manager ist noch nicht ganz abgeschlossen, allerdings stehen schon grobe Seitenstrukturen und erste Entwürfe vom Design einzelner Seiten. 

In der Konzeption gibt es bisher eine Startseite, auf der man eine Hero Section mit einführenden Informationen hat. Weiter darunter ist dann die Eventsuche und einige Events kategorisiert in \enquote{Neue Events} und Events, die demnächst anstehen. Außerdem gibt es auch noch eine Auflistung aller Events.

Über die Suche kann man Events nach Name und Ort suchen. Sobald die Suche ausgeführt wird, wird man auf eine Unterseite mit den Suchergebnissen weitergeleitet. Zusätzlich kann man über die angezeigten Events auch eine Detail-Seite zu jedem Event aufrufen. Dort werden dann detailliertere Informationen, wie Detail-Beschreibung, Google Maps usw., angezeigt. Außerdem hat man hier die Möglichkeit das Event zu buchen. Zu den eben genannten Seiten existierte bereits ein Design. Zu den folgenden genannten Seiten existieren bisher nur LoFi-Schemas, um zu zeigen, was der Inhalt und ungefähre Aufbau der Seite ist.

Außerdem soll es möglich sein, ein Event zu erstellen. Um weitere Formulare in dieser Anwendung barrierefrei umzusetzen, ist auch noch eine Kontaktseite mit Kontaktformular, Login und Registrierung geplant. Dazu soll man auch ein Event buchen können. Die Seiten zielen darauf ab, Konzepte der Barrierefreiheit anhand passender Komponenten zu erklären.

\section{Design und Konzeption}
Da noch nicht alle Seiten im Design umgesetzt wurden und die Seiten sich noch in der Entwicklung befinden, wurde für diese Arbeit ein eigenes Designsystem entwickelt, um den Fokus auf die Barrierefreiheit legen zu können. Dabei orientiert sich das Designsystem aber stark an Mrija Manager.

Daher weicht die Seitengestaltung auch von der Vorgabe ab. Damit bietet sich die Möglichkeit, bei bestimmten Elementen die Grenzen der Barrierefreiheit zu dokumentieren.  So wurde zum Beispiel noch ein FAQ hinzugefügt, das mit einem Akkordion System umgesetzt wurde. Anhand der Interaktivität soll die Umsetzung gerade im Bereich der Barrierefreiheit gezeigt werden.
\subsection{Farbkonzept}
Das Farbkonzept wurde vollständig aus den bisherigen Planungen übernommen, um zu zeigen, wie man Farben nutzen kann, bei denen bei der Auswahl Barrierefreiheit nicht berücksichtigt wurde. Die Farbpalette wurde deshalb im ersten Schritt im Contrast Grid auf Kontraste geprüft. Es gibt auch einen offiziellen W3c Contrast Checker, allerdings ist der Vorteil vom Contrast Grid, dass alle möglichen Farbkombinationen auf den ersten Blick erkennbar sind. Das Contrast Grid mit der verwendeten Farbpallette ist hier abrufbar: \href{https://contrast-grid.eightshapes.com/?version=1.1.0&background-colors=&foreground-colors=%231C1B1F%2C%20Dark%0D%0A%23000000%2C%20Black%0D%0A%23FFFFFF%2C%20White%0D%0A%23666666%2C%20Black%2060%25%2C%0D%0A%23C4C4C4%2C%20Black%2023%25%2C%0D%0A%23757575%2C%20Black%2054%25%2C%0D%0A%23949494%2C%20Black%2048%25%2C%0D%0A%23F3FAFE%2C%20Light%20Blue%20%2B1%0D%0A%23E2F0F4%2C%20Light%20Blue%0D%0A%2384234F%2C%20Accent%0D%0A%23F5F5FF%2C%20Primary%20%2B2%0D%0A%239999C5%2C%20Primary%20%2B1%0D%0A%234748C6%2C%20Primary%0D%0A%232E2F5C%2C%20Primary%20-1%0D%0A%23000450%2C%20Primary%20-2&es-color-form__tile-size=compact&es-color-form__show-contrast=aaa&es-color-form__show-contrast=aa}{Contrast Grid}\footnote{\url{https://www.figma.com/proto/jnGNE7gzjqBn1uD3rOXifO/bachelor-thesis?node-id=1-79\&p=f\&t=NlQQAOGGDPRokds2-0\&scaling=min-zoom\&content-scaling=fixed\&page-id=0\%3A1\&starting-point-node-id=1\%3A79}}

Auf die genaue Gestaltung der Seite wird hier nicht weiter eingegangen, denn die Konzepte der Barrierefreiheit lassen sich in der Programmierung besser erklären und der Fokus liegt hier auf der Programmierung.
